\newpage
\begin{center}
  \textbf{\large 2 Описание метода}
\end{center}
\refstepcounter{chapter}
\addcontentsline{toc}{chapter}{2 Описание метода}

\section{Выбор базового метода}
На основе данных проведенного обзора можно сделать вывод о наличии тренда стремительного развития
подходов с применением глубокого обучения в задаче одометрии не смотря на превосходство
классических методов в точности оценок. В частности, использование
глубоких нейронных сетей в задачах объединения данных разных модальностей позволяет
получить качественные обобщения, что может быть применено и в смежных областях. Кроме
того, неявные обобщения могут быть значительно более эффективными.

Метод DVLO~\cite{liu2024dvlo} показывает себя как один из наиболее качественных. Его точность
превосходит большинство нейросетевых мультимодальных методов, а также часть классических. Заявленные
результаты приведены в таблице~\ref{comp-Table1}.

\begin{table}[h]
  \caption{Сравнение подходов одометрии, на данных KITTI Odometry Benchmark~\cite{geiger2013vision}. $t_{\mathrm{rel}}$ и $r_{\mathrm{rel}}$ --- относительная ошибка расстояния и относительная ошибка ориентации на 100 м соответственно}
  \label{comp-Table1}
  \centering{
    \begin{tabular}{ |C{4cm}|C{4cm}|C{1.5cm}|C{3cm}| }
      \hline
      Метод                                & Модальность    & $t_{\mathrm{rel}}, \%$ & $r_{\mathrm{rel}},$ $^\circ/100$м \\ \hline
      \multicolumn{4}{|c|}{\textit{Нейросетевые методы}}                                               \\
      \hline
      DVLO~\cite{liu2024dvlo}              & LiDAR + камера & $0.82$        & $0.41$                   \\
      H-VLO~\cite{aydemir2022h}            & LiDAR + камера & $1.36$        & $0.43$                   \\
      EfficientLO~\cite{wang2022efficient} & LiDAR          & $0.86$        & $0.41$                   \\
      TransLO~\cite{liu2023translo}        & LiDAR          & $0.99$        & $0.50$                   \\
      D3VO~\cite{yang2020d3vo}             & камера         & $0.88$        & $0.21$                   \\
      DVSO~\cite{yang2018deep}             & камера         & $0.9$         & $0.21$                   \\
      \hline
      \multicolumn{4}{|c|}{\textit{Классические методы}}                                               \\
      \hline
      VLOAM~\cite{7139486}                 & LiDAR + камера & $0.75$        & ---                      \\
      DVL-SLAM~\cite{shin2020dvl}          & LiDAR + камера & $0.98$        & ---                      \\
      PL-LOAM~\cite{9196613}               & LiDAR + камера & $0.94$        & $0.36$                   \\
      KISS-ICP~\cite{vizzo2023kiss}        & LiDAR          & $0.61$        & $0.17$                   \\
      SOFT2~\cite{cvivsic2022soft2}        & камера         & $0.53$        & $0.1$                    \\
      \hline
    \end{tabular}
  }
  
\end{table}

Рассматриваемый метод представляет уникальный подход к объединению данных, исследование которого
может способствовать развитию данного подхода. В связи с этим в качестве базового был выбран
именно DVLO~\cite{liu2024dvlo}.

\section{Принцип работы}

\subsection{Изображение как набор точек}
\label{sec:set_of_points}
Важную роль в рассматриваемом методе играет особый способ представления изображений
как набора точек на основе контекстной кластеризации (CoC)~\cite{ma2023image}. Идея заключается
в использовании нового подхода для экстракции признаков при работе с визуальными данными,
позволяющем увеличить эффективность при качестве сравнимом с широко применяемыми
сверточными нейронными сетями (CNN)~\cite{simonyan2014very} и визуальными трансформерами (ViT)~\cite{dosovitskiy2020image}.

Конвенционально изображения представлялись в качестве прямоугольного $h\times w$ трехканального массива
$I \in \mathbb{R}^{h\times w \times 3}$. В данном подходе предлагается использовать пространственную
информацию о положении пикселя в качестве непосредственно первичного признака, а все изображение
рассматривать как неупорядоченный набор точек с обобщенными признаками $I^P \in \mathbb{R}^{n_p \times 5}$,
где $n_p = h\cdot w$. При этом пространственная информация каждого пикселя $I_{i,j}$
представляется следующим образом: $\begin{bmatrix} \frac{i}{h}-0.5 & \frac{j}{w}-0.5 \end{bmatrix}$.

После преобразования изображений производится иерархическая экстракция признаков. Первым шагом
является редуцирование набора точек посредством равномерного выбора ключевых точек и
проецирования некоторой окрестности при помощи линейного слоя в единое представление.
После редуцирования применяется процедура контекстной кластеризации, что и является ключевой
идеей предлагаемого метода. Редуцированные точки проецируются в пространство сравнения при помощи
полносвязной нейронной сети. Среди полученных точек $I^p_s \in \mathbb{R}^{n_r \times d}$
(количество признаков может меняться при проецировании) также равномерно выбираются
центры будущих $c$ кластеров --- их признаки рассчитываются как среднее среди $k$ ближайших соседей.
Затем рассчитывается карта сходства $S \in \mathbb{R}^{c\times n_r}$ на основе
косинусного расстояния между центрами кластеров и точками. Таким образом,
используя карту сходства можем определить принадлежность точек к кластерам.
Имея кластер, состоящий из $m$ точек и соответствующие им значения сходства $s \in \mathbb{R}^m$
производится проецирование в пространство значений $I^p_v \in \mathbb{R}^{n_r\times d'}$ (также проецируется центр кластера $v_c$).
Рассчитывается агрегированный признак класса:
\begin{equation*}
  g = \frac{1}{C} \left( v_c + \sum_{i=1}^m \sigma(a s_i + \beta) \cdot v_i \right), \quad  C = 1 + \sum_{i=1}^m \sigma(a s_i + \beta),
\end{equation*}
где $v_i$ --- значение для $i$-й точки кластера, $\sigma$ --- сигмоидальная функция, $\alpha$ и $\beta$ --- обучаемые параметры.
После получения $g$ предлагается обновить признаки каждой точки кластера. Пусть $p_i \in \mathbb{R}^{d}$ ---
исходный вектор признаков $i$-й точки. Итоговое обновление выполняется формулой~\eqref{eq:coc_update}
\begin{equation}
  \label{eq:coc_update}
  p'_i \;=\; p_i \;+\;
  \mathrm{FC}\Bigl(\sigma(\alpha\, s_i + \beta)\,\cdot\, g\Bigr),
\end{equation}
где $\mathrm{FC}$ --- полносвязный слой, отображающий из $\mathbb{R}^{d'}$ в $\mathbb{R}^{d}$.
Таким образом, каждая точка получает <<распространённый>> признак, пропорциональный величине сходства $s_i$.
Данная процедура повторяется несколько раз в зависимости от количества блоков в иерархическом
модуле.

\subsection{Иерархическая экстракция признаков}
Важную роль в представлении облаков точек, полученных с лидара ${P}_{\mathrm{lidar}} \;=\; \bigl\{\,{p}_i \in \mathbb{R}^3 \mid i=1,\ldots,N\bigr\}$
играет переход к псевдоизображению. Для этого используется цилиндрическая проекция, которая задаётся формулами~\eqref{eq:cylind_proj}.
\begin{equation}
  \label{eq:cylind_proj}
  u \;=\; \frac{1}{\Delta \theta}\,\arctan2(y,\,x),
  \qquad
  v \;=\; \frac{1}{\Delta \varphi}\,\arcsin\!\Bigl(\frac{z}{\sqrt{x^2 + y^2 + z^2}}\Bigr),
\end{equation}
где $(x,y,z)$~--- координаты точки в системе лидара, а
$\Delta \theta,\ \Delta \varphi$~--- шаг углов по горизонтали и вертикали соответственно.
Таким образом, каждая точка проецируется в пиксель $(u,v)$.
В соответствующую позицию записываются значения координат в результате чего формируется тензор (далее --- псевдоизображение):
\begin{equation*}
  {I}_{\mathrm{lidar}} \;\in\; \mathbb{R}^{H_P \times W_P \times 3}.
\end{equation*}

Для входного изображения ${I}_{\mathrm{cam}} \;\in\; \mathbb{R}^{H \times W \times 3}$
применяется свёрточная сеть для выделения признаков в виде ${F}_I \;\in\; \mathbb{R}^{H_{I} \times W_{I} \times C_{I}},$
где $H_{I}, W_{I}$ могут быть меньше исходных $H, W$, а $C_{I}$~--- количество каналов признаков.
Аналогичным образом извлекаются признаки из псевдоизображений: 
${F}_P \;\in\; \mathbb{R}^{H_P \times W_P \times C_{p}}$.

Описанные выше операции производятся несколько раз последовательно ---
в оригинальной архитектуре предусмотрено 4 уровня экстракции.


\subsection{Локальное объединение признаков (Local Fuser)}
\label{sec:local_fuser}
Предложенный механизм локального объединения (Local Fuser), позволяет объединять текстурные 
признаки изображения и геометрические признаки лидара в локальной окрестности. 
Основная идея состоит в представлении карты признаков изображения как множества
псевдоточек аналогично \ref{sec:set_of_points}, а каждый элемент облака точек при 
этом рассматривается в качестве центра контекстного кластера в дальнейшей обработке.

Пусть имеется карта признаков изображения $F_{I} \;\in\; \mathbb{R}^{H_{I} \times W_{I} \times C_I}$.
Путём преобразования размерности формируется набор псевдоточек $F_{pp} \in \mathbb{R}^{M\times C_I}$, где $M=H_I\times W_I$
и каждая из них отвечает отдельной позиции в $F_{I}$.

Пусть лидарное облако ${P}_{\mathrm{lidar}}$ состоит из $N$~3D-точек. При известной калибровочной матрице 
${K}[{R}|{t}]$ можно спроецировать каждую ${p}_i \in {P}_{\mathrm{lidar}}$ в плоскость изображения согласно уравнению~\eqref{eq:l2cam_proj}. 
\begin{equation}
  \label{eq:l2cam_proj}
  (x'_i,\, y'_i) \;=\; \Pi_{\mathrm{cam}}\!\bigl({p}_i\bigr),
\end{equation}
где $\Pi_{\mathrm{cam}}$ обозначает процедуру проекции в координаты изображения.
Затем при помощи билинейной интерполяции извлекаются соответствующие признаки из ${F}^{I}$:
\begin{equation*}
  {c}_i \;=\; \mathrm{Interp}\!\Bigl({F}_{I},\, x'_i,\, y'_i\Bigr)
  \;\in\;\mathbb{R}^{C_I}.
\end{equation*}
Каждый такой вектор ${c}_i$ интерпретируется как центр кластера. 
Таким образом, общее число центров $N$ совпадает с числом лидарных точек.

Пусть \[
\mathrm{Sim}\!\bigl(i,j\bigr) 
  \;=\;
  \frac{{c}_i^\top\,{p}_j}
       {\|{c}_i\|\,\|{p}_j\|} 
  \quad\in\;\mathbb{R}
\]
обозначает косинусное сходство между $i$-м центром ${c}_i$ и $j$-й псевдоточкой ${p}_j$ из $F_{pp}$. 
Тогда для $i$-го кластера вычисляется агрегированный признак согласно формуле~\eqref{eq:local_fuse_agg}.
\begin{equation}
\label{eq:local_fuse_agg}
  {g}_i
  \;=\;
  \frac{1}{1 \;+\; \sum_{j=1}^M\,\sigma\bigl(\alpha\, \mathrm{Sim}(i,j) \;+\;\beta\bigr)}\;
  \Bigl(\,
    {c}_i 
    \;+\;
    \sum_{j=1}^M\,\sigma\bigl(\alpha\, \mathrm{Sim}(i,j) \;+\;\beta\bigr)\,{p}_j
  \Bigr),
\end{equation}
где $\alpha,\beta$ --- обучаемые скалярные параметры, а $\sigma(\cdot)$ --- сигмоидальная функция. 

Вектор ${g}_i \in \mathbb{R}^{C_I}$ становится {локально-объединенным} признаком для $i$-й лидарной точки. 
Итоговый тензор локальных признаков $F_L \in \mathbb{R}^{H_P\times W_P \times C_I}$ составлен из $g_i$ на 
основе данных о соответствии проецированных точек.

\subsection{Глобальное объединение признаков (Global Fuser)}
\label{sec:global_fuser}

После выполнения локального объединения, каждая лидарная точка обладает агрегированным признаком 
${g}_i \in \mathbb{R}^{C_I}$, учитывающим локальную текстурную информацию из карты камеры. 
Однако локальное объединение охватывает лишь окрестности в признаковом пространстве, 
не обеспечивая достаточной глобальной взаимосвязи. 
Чтобы учесть более глобальный контекст и далее согласовать лидарные и визуальные особенности, 
используется выход иерархического экстрактора признаков и осуществляется адаптивное объединение на уровне всей карты.

В итоге на вход поступает две карты признаков: 
\[
  {F}_P\in\mathbb{R}^{H_P\times W_P\times C_P}
  \quad\text{и}\quad
  {F}_L\in\mathbb{R}^{H_P\times W_P\times C_I},
\]
где ${F}_P$ --- признаки облака точек лидара, 
а ${F}_L$ --- результат локального объединения, учитывающий текстуры изображения. 
Для их глобального объединения вводится набор обучаемых весов, 
которые для каждой позиции $(u,v)$ регулируют вклад ${F}_P$ и ${F}_L$. 
В частности, можно ввести адаптивные карты ${A}_P,\,{A}_L$ (одинакового размера), вычисляя их через небольшие полносвязные (MLP) или свёрточные слои:
\begin{equation*}
\label{eq:global_adaptive_weights}
  {A}_P(u,v) \;=\; \sigma\!\Bigl(\mathrm{MLP}\bigl({F}_P(u,v)\bigr)\Bigr), 
  \quad
  {A}_L(u,v) \;=\; \sigma\!\Bigl(\mathrm{MLP}\bigl({F}_L(u,v)\bigr)\Bigr),
\end{equation*}
где $\sigma(\cdot)$ --- сигмоидальная функция, сжимающая выход в интервал $(0,1)$. 
Тогда итоговый признак для позиции $(u,v)$:
\begin{equation}
\label{eq:global_fused_feature}
  {F}_G(u,v)
  \;=\; 
  \frac{
    {A}_P(u,v)\,\odot\,{F}_P(u,v)
    \;+\;
    {A}_L(u,v)\,\odot\,{F}_L(u,v)
  }{
    {A}_P(u,v) \;+\;{A}_L(u,v)\;+\;\varepsilon
  },
\end{equation}
где символ $\odot$ обозначает покомпонентное умножение, 
а $\varepsilon$ --- малая константа для числовой стабильности.

Полученная карта
\(
  {F}_G 
  \;\in\; 
  \mathbb{R}^{H_P\times W_P\times d}
\)
содержит как исходную лидарную информацию, 
так и локально-агрегированные признаки, 
взвешенные согласно обучаемым коэффициентам в~\eqref{eq:global_fused_feature},
обеспечивая глобальный контекст и тем самым повышая устойчивость к шумам и неоднозначностям, 
связанным с локальными несоответствиями.

\subsection{Итеративная оценка положения (Odometry Regression)}
\label{sec:iter_pose_est}

После получения глобально объединенных признаков ${F}_G^S,\;{F}_G^T$
для двух последовательных кадров (источника --- <<$S$>> и цели --- <<$T$>>) 
необходимо определить относительное перемещение между ними, то есть 
кватернион вращения ${q}\in\mathbb{R}^4$ и вектор смещения ${t}\in\mathbb{R}^3$.

Признаки ${F}_G^S$ и ${F}_G^T$ могут быть далее ассоциированы через механизм вычисления тензора соответствий 
(cost embedding), отражающего, насколько каждая позиция в кадре~$S$ соответствует аналогичной позиции в кадре~$T$. 
Можем определить:
\[
  {E}
  \;=\;
  \mathrm{CostEmbed}\Bigl({F}_G^S,\;{F}_G^T\Bigr)
  \;\in\;\mathbb{R}^{H_p \times W_p \times {d}},
\]
где $\mathrm{CostEmbed}(\cdot,\cdot)$ --- блок свёрточных слоёв и MLP, вычисляющий сходство между парами признаков. 
${E}$ представляется в виде массива $\mathbb{R}^{N\times d}$. В дальнейшем $F_G$ приведены к такой же размерности.

Чтобы извлечь из ${E}$ итоговые $({q},{t})$, вводится небольшой регрессионный блок. 
Вычисляется карта ${M}\in\mathbb{R}^{H_p \times W_p \times d}$, которая содержит информацию о значимых областях,
учитывая признаки ${E}$ и ${F}_G^S$, после чего суммируется по пространственным координатам:
\[
  {M} \;=\; \mathrm{softmax}\Bigl(\mathrm{MLP}\bigl({E}\,\oplus\,{F}_G^S\bigr)\Bigr),
  \quad
  {m} \;=\;
    \sum_{u,v}\,
      {M}(u,v)\,\odot\,{E}(u,v),
\]
где $\oplus$ означает конкатенацию по канальной оси. В результате ${m}\in\mathbb{R}^{d}$ (после суммирования или свёрток) служит входом для финального линейного слоя~\eqref{eq:final_qt}.
\begin{equation}
  \label{eq:final_qt}
  \bigl(\Delta {q}, \Delta {t}\bigr)
  =
  \mathrm{FC}\bigl({m}\bigr),
\end{equation}
где $\Delta {q}\in\mathbb{R}^4$ и $\Delta {t}\in\mathbb{R}^3$ --- предсказанные приращения, которые формируют текущую оценку $({q},{t})$:
\begin{equation}
\label{eq:iter_q_update}
  {q}_{l} 
  \;=\; 
  \Delta {q}_{l}\,\otimes\,{q}_{l+1}, 
  \qquad
  {t}_{l}
  \;=\; 
  {t}_{l+1} 
  \;+\; \mathrm{Rot}\bigl(\Delta {q}_{l}\bigr)\,\Delta {t}_{l},
\end{equation}
где $\otimes$ обозначает умножение кватернионов, а $\mathrm{Rot}(\Delta {q})$ --- матрица поворота, полученная из кватерниона $\Delta {q}$. 
Таким образом, на каждом уровне $l$ уточняется грубая оценка $({q}_{l+1}, {t}_{l+1})$ 
и в конце получаем итоговое $({q}_1, {t}_1)$.

Чтобы модель могла извлечь грубые соответствия на низком разрешении и затем уточнить их с учётом более <<детальных>> признаков, 
процедура~\eqref{eq:iter_q_update} повторяется на каждом уровне экстракции и объединения признаков. 
При этом формируется свой собственный cost embedding ${E}_l$ и предсказываются $(\Delta {q}_l, \Delta {t}_l)$. 
Наконец, итоговая предсказанная поза $({q}_1, {t}_1)$ имеет существенно меньшую ошибку, чем грубая начальная $({q}_4, {t}_4)$.

\subsection{Функция потерь}
\label{sec:loss_function}

Обучение сети происходит в парадигме обучения с учителем,
опирающейся на эталонные относительные позы.
Соответствующие опорные значения обозначаются следующим образом
${q}_{\mathrm{gt}}$ и ${t}_{\mathrm{gt}}$.

Для каждого из четырёх иерархических уровней вводится частичная
функция потерь:
\begin{equation}
\label{eq:level_loss}
\mathcal{L}_{l} \;=\;
\bigl\lVert {t}_{\mathrm{gt}} - {t}_{l} \bigr\rVert_{1}\,
e^{-k_{x}}
\;+\;
k_{x}
\;+\;
\bigl\lVert {q}_{\mathrm{gt}} - {q}_{l} \bigr\rVert_{2}\,
e^{-k_{q}}
\;+\;
k_{q},
\end{equation}
где 
$k_{x},\,k_{q}\in\mathbb{R}$ — обучаемые скалярные параметры,
служащие для автоматического подбора весов между
трансляционной и вращательной составляющими ошибки.
Кватернионы нормируются к единичной длине перед
вычислением нормы.

Совокупная функция потерь формируется как взвешенная сумма частичных функций из уравнения~\eqref{eq:level_loss}:
\begin{equation*}
\label{eq:total_loss}
\mathcal{L} \;=\;
\sum_{l=1}^{4} \alpha^{l}\,\mathcal{L}_{l},
\end{equation*}
где коэффициенты $\alpha^{l}$ отражают вклад
каждого уровня в итоговую оптимизацию.

Наличие экспоненциальных весов $e^{-k_{x}}$ и
$e^{-k_{q}}$ обеспечивает адаптивное согласование
масштабов двух слагаемых, что избавляет от ручного
подбора статических коэффициентов. Обратные
добавления $k_{x}$ и $k_{q}$ предотвращают неограниченный
рост соответствующих параметров, сохраняя градиенты
ненулевыми даже при малых ошибках.

\section{Модификации}
\subsection{Стерео-камера}
\label{sec:stereo_camera}

В качестве одной из модификаций для улучшения качества работы модели 
предлагается интегрировать в архитектуру сети возможность использовать 
стерео-камеру: использовать два изображения --- левое и правое, 
каждое из которых соответствует собственному ракурсу. 
Дополнительная информация о сцене, получаемая за счёт стереопары 
может повысить робастность и точность итоговой одометрии, 
особенно при наличии частичных перекрытий или динамических объектов.

В исходном подходе DVLO для этапа локального объединения
каждая лидарная точка $p_i \in \mathbb{R}^3$ выступает центром кластера, 
проецируясь на одно изображениею
В предлагаемом расширении предполагается, что 
каждая лидарная точка проецируется на оба изображения. 
Далее эти признаки конкатенируются, для получения более полной локальной картины.

Пусть ${F}_{I_L} \in \mathbb{R}^{H_I \times W_I \times C_I}$ и ${F}_{I_R} \in \mathbb{R}^{H_I \times W_I \times С_I}$
--- карты признаков для левого и правого изображений соответственно.
Для локального объединения с камерой, как и прежде в формуле~\eqref{eq:l2cam_proj}, формируются псевдоточки из каждого изображения: $F_{pp_L}$ и $F_{pp_R}$.
Затем каждая лидарная точка ${p}_i$ проецируется как 
\begin{equation*}
  (x'_i,\,y'_i)_L = \Pi_L({p}_i), 
  \quad
  (x'_i,\,y'_i)_R = \Pi_R({p}_i),
\end{equation*}
где $\Pi_L,\,\Pi_R$ --- функции проекции с учётом калибровок соответствующих камер. 
Из ${F}_{I_L}$ и ${F}_{I_R}$ извлекаются центровые признаки ${c}_i^L,\,{c}_i^R \in \mathbb{R}^{C}$ 
(посредством билинейной интерполяции), которые играют роль кластерных центров для локального объединения.

Предлагается следующий вариант конкатенации:
\begin{equation*}
  {c}_i^{\mathrm{concat}} 
  \;=\;
  \mathrm{Concat}\bigl({c}_i^L,\;{c}_i^R \bigr)
  \;\in\;\mathbb{R}^{2C_I},
\end{equation*}
то есть центры из левого и правого изображения объединяются в один вектор. 
Аналогично псевдоточки конкатенируются в общий тензор:
\[
  F_{pp}^{\mathrm{concat}} 
  \;=\;
  F_{pp_L} \;\cup\; F_{pp_R},
\]
и в процессе локального объединения на основе косинусного сходства 
все псевдоточки распределяются по классам в соответсвии с близостью к ${c}_i^{\mathrm{concat}}$. 
Таким образом, каждая лидарная точка получает обогащённый признак, 
который учитывает информацию из обеих камер одновременно.

Аналогично исходному методу, после получения локально объединенных признаков ${F}_L$ происходит расчет глобальных признаков
с помощью адаптивного механизма~\eqref{eq:global_fused_feature}. 
Разница заключается лишь в том, что теперь ${F}_L$ внутри уже содержит стерео-информацию. 